\section{Compact factors, multiple averages, and Haar rigidity}
\label{sec:rigidity}

We complete the proof of Theorem~\ref{thm:rigidity}.  The first step
constructs a factor to which Lemma~\ref{lem:relative-singularity} applies
and above which conditional correlations decouple.  Two identities among
algebraic unit directions then turn this decoupling into a difference
inequality for a radial function.  A finite number-field argument forces
that function to vanish away from zero.

Throughout this section a relative product uses conditional independence
over the indicated factor.  If $\pi:X\to Y$ is a factor map and
$\mu=\int\mu_y\,d\nu(y)$ is its disintegration, its relative product has
measure
\[
 \mu\mathbin{\times_Y}\mu
   =\int \mu_y\otimes\mu_y\,d\nu(y)
\]
on $X\times_YX$.  The action on this product is diagonal.  All factors
are understood modulo null sets; countability of $K$ permits all the
countably many equivariance identities to be imposed simultaneously.

\subsection{A compact-factor tower}

The following construction is the compact/relative-weak-mixing part of
the Furstenberg--Zimmer structure theorem and relative dichotomy
\citep{Furstenberg1977,ZimmerExtensions1976,ZimmerDiscrete1976,
JamneshanStructure,JamneshanCorrigendum}. The passage from invariant
Hilbert--Schmidt kernels to finite-rank modules, and the boundedness
of their bases in the ergodic case, appear in
\citet[Lemmas~6.6 and~7.2]{Furstenberg1977}. We give the fiberwise
proof because the particular approximation property
\eqref{eq:source-4} is needed here.

\begin{lemma}[Compact tower]\label{lem:compact-tower}
Let $X$ be an ergodic standard probability-preserving $K$-system.
There is a countable ordinal $\eta$ and an increasing tower of factors
$(Y_\alpha)_{\alpha\leq\eta}$ with the following properties:
\begin{enumerate}
\item $Y_0$ is the trivial factor, and each limit stage is the factor
generated by its predecessors;
\item every successor extension $Y_{\alpha+1}\to Y_\alpha$ satisfies
\eqref{eq:source-4};
\item for $Y=Y_\eta$, the system $X\times_YX$ is ergodic.
\end{enumerate}
\end{lemma}

\begin{proof}
We first construct a strict extension $Y'$ of any factor $Y$ whose
relative product is not ergodic.  An invariant kernel on that product
will provide an equivariant finite-rank space on almost every fiber.  Its
bounded basis generates $Y'$ and supplies the approximation property
\eqref{eq:source-4}.  We then iterate this enlargement.

Suppose that $Y$ has already been constructed and that $X\times_YX$ is
not ergodic.  Choose a bounded invariant nonconstant kernel $L(x,x')$
on this relative product.  Its conditional expectations onto either
coordinate are invariant functions on $X$, so both are the constant
$\int L\,d(\mu\mathbin{\times_Y}\mu)$.  Subtract this constant.  The
result is still nonzero, and its two coordinate conditional expectations
are zero.

For almost every $y$, let $A_y$ be the integral operator with kernel
$L$ on the fiber Hilbert space $\mathcal H_y=L^2(\mu_y)$.  These are
Hilbert--Schmidt operators, with
\[
 \norm{A_y}_{\mathrm{HS}}\leq\norm{L}_\infty,
 \qquad A_y1=A_y^*1=0.
\]
Invariance of the kernel makes the operator field equivariant under the
unitary maps between fibers.  For a sufficiently small fixed $\delta>0$,
the spectral projection
\[
 Q_y=\ind_{(\delta,\infty)}(A_y^*A_y)
\]
is nonzero on a positive-measure set of fibers.  Its rank is finite and
bounded above by $\norm{L}_\infty^2/\delta$.  The rank is invariant on
$Y$, and $Y$ is ergodic as a factor of $X$.  It is therefore a constant
integer $d\geq1$ almost everywhere.  The range of $Q_y$ is orthogonal
to the fiber constants.

Here are the measurability details for this construction.  Work with
standard Borel models and fiber-supported conditional probabilities.
Choose a countable algebra generating the Borel structure of $X$ and
its simple functions with complex-rational coefficients.  Their
restrictions are dense in each $L^2(\mu_y)$.  Conditional integration
against the bounded kernel represents the images of these sections under
$A_y$ and $A_y^*$ by measurable functions on $X$; the same holds for
polynomials in $A_y^*A_y$.  On their common bounded spectral interval,
there are uniformly bounded polynomials converging pointwise to
$\ind_{(\delta,\infty)}$: first choose bounded continuous approximations
to this indicator, with value zero at $\delta$, and then uniformly
approximate them by polynomials.  Spectral calculus gives convergence
in each fiber Hilbert norm.  Domination gives convergence in global
$L^2(X)$ as well.  A subsequence converges in fiber norm for almost every
$y$, identifying the global limit with the projected section.  Thus the
projected dense sections are measurable.  Their fiber scalar products
make rank and linear-independence tests measurable.  Selecting the first
available independent section at each step and applying fiberwise
Gram--Schmidt gives measurable fiber-orthonormal sections
$e_1,\ldots,e_d$ spanning the range of $Q_y$.

For each $k\in K$ the sections $U_ke_i$ differ from the $e_i$ by a
unitary coefficient matrix over $Y$.  Consequently
$\sum_{i=1}^d|e_i(x)|^2$ is invariant on $X$.  Its integral is $d$, so
ergodicity gives
\begin{equation}\label{eq:rigidity-bounded-basis}
 \sum_{i=1}^d|e_i(x)|^2=d\quad\text{almost everywhere}.
\end{equation}
In particular, the sections are bounded.  Adjoin them to $Y$ to generate
an invariant factor $Y'$.  This is a strict extension: every $e_i$ has
conditional mean zero over $Y$ and has unit fiber norm, and hence is not
$Y$-measurable.

Bounded polynomials in the $e_i,\bar e_i$, with bounded coefficients
from $Y$, are dense in $L^2(Y')$.  Indeed these functions form a bounded
unital conjugation-closed algebra generating $Y'$, and approximation
first by functions of finitely many generators and then by polynomials
gives the density.  A rotation of a fixed polynomial expands into
finitely many fixed bounded monomials in the $e_i,\bar e_i$.  Its
coefficients are functions on $Y$, uniformly bounded independently of
the rotation: the original coefficient sup norms are preserved, and
the entries of the unitary matrices have modulus at most one.  For a
bounded $H$ on $Y'$, choose one polynomial approximation in $L^2$.
Rotating preserves its approximation error.  The finite expansion of
the rotated polynomial proves \eqref{eq:source-4} for $Y'\to Y$.

Starting at the trivial factor, repeat this step whenever the relative
product is not ergodic, and take generated joins at limit ordinals.
Every countable stage is a standard factor modulo null sets.  The
process stops at a countable ordinal.  Otherwise, at every successor
through all countable ordinals, choose a unit vector in the new factor
orthogonal to the preceding factor.  Nestedness makes these vectors
pairwise orthogonal. There would be $\aleph_1$ such vectors,
contradicting separability of $L^2(X)$.  At the
stopping stage the relative product is ergodic by the stopping rule.
\end{proof}

\subsection{Conditional multiple averages}

A F\o lner sequence in $K$ means nonempty finite sets $\Phi_N\subset K$
such that
\[
 \frac{|h\Phi_N\mathbin{\triangle}\Phi_N|}{|\Phi_N|}\longrightarrow0
 \qquad(h\in K).
\]
Such sequences exist because $K$ is countable abelian.  Write
$\operatorname{Avg}_{u\in\Phi}b_u=|\Phi|^{-1}\sum_{u\in\Phi}b_u$.
Every average below is a finite-set average in this discrete group.

We use the Hilbert-space mean ergodic theorem in its every-F\o lner
form.  For a unitary representation $V$ of $K$, these averages converge
strongly to the invariant projection along every F\o lner sequence.
For completeness, the orthogonal complement of the invariant vectors
is the closed span of vectors $V_hv-v$.  The average of such a vector
has norm at most
$|h\Phi_N\triangle\Phi_N|\norm{v}/|\Phi_N|$, and density proves the
claim.  In particular, the representation $u\mapsto U_{u^d}$, for
$d\in\Z\setminus\{0\}$, has the same invariant vectors as $U$, because
the power map $u\mapsto u^d$ is onto.  No assertion about unweighted
images of F\o lner sets under a power map is needed.

The proof below uses the relative weak-mixing induction for multiple
averages, with a Hilbert-space van der Corput estimate; compare
\citet{Furstenberg1977}.  We give the argument for the power maps of
$K$ along every F\o lner sequence, since both features are used below.

\begin{lemma}[Conditional multiple averages]\label{lem:multiple-averages}
Suppose that $X$ is an ergodic standard probability-preserving
$K$-system, that $Y$ is an invariant factor, and that $X\times_YX$ is
ergodic.  Let $P=\mathbb E(\,\cdot\mid Y)$.  For pairwise distinct
integers $n_1,\ldots,n_l$ and bounded functions $v_1,\ldots,v_l$ on $X$,
\begin{equation}\label{eq:source-6}
 \operatorname{Avg}_{u\in\Phi_N}
 \left\|P\left(\prod_{j=1}^lU_{u^{n_j}}v_j\right)
       -\prod_{j=1}^lU_{u^{n_j}}Pv_j\right\|_2^2
 \longrightarrow0
\end{equation}
along every F\o lner sequence $(\Phi_N)$ in $K$.  The pair criterion
used in the proof is, for bounded $v,w$ and $d\ne0$,
\begin{equation}\label{eq:source-7}
 \operatorname{Avg}_{u\in\Phi_N}
 \norm{P(vU_{u^d}w)-Pv\,U_{u^d}Pw}_2^2\longrightarrow0.
\end{equation}
\end{lemma}

\begin{proof}
Put $\mathcal R=X\times_YX$ and denote its conditional expectation onto
$Y$ by $P_{\mathcal R}$.  First prove \eqref{eq:source-7}.  With
$v_0=v-Pv$ and $w_0=w-Pw$, its difference equals
$P(v_0U_{u^d}w_0)$.  Conditional independence gives
\begin{equation}\label{eq:rigidity-pair-tensor}
 \norm{P(v_0U_{u^d}w_0)}_2^2
 =\int_{\mathcal R}(v_0\otimes\bar v_0)
       U_{u^d}(w_0\otimes\bar w_0)\,d\mu_{\mathcal R}.
\end{equation}
The integral of $w_0\otimes\bar w_0$ is
$\int_Y|Pw_0|^2\,d\nu=0$.  The mean ergodic theorem on the ergodic
system $\mathcal R$, applied through the surjective power map, makes
the average of the right side tend to zero.  This proves the pair
criterion along every F\o lner sequence.

The same criterion holds on $\mathcal R$ over $Y$.  To see this first
for simple tensors $v=a_1\otimes a_2$ and $w=b_1\otimes b_2$, use
\[
 P_{\mathcal R}(vU_{u^d}w)
 =P(a_1U_{u^d}b_1)P(a_2U_{u^d}b_2).
\]
Subtract the analogous product of the factor parts.  The difference is
a sum of two terms, each containing one error from
\eqref{eq:source-7} on $X$ and one uniformly bounded factor.  Its
averaged squared $L^2(Y)$ norm therefore tends to zero.  Linearity gives
the result for finite sums of tensors.  Such sums can approximate any
bounded function on $\mathcal R$ in $L^2$, with a common sup bound:
use conditional expectations onto finite partitions generated by
rectangles in the two coordinates.  Approximate one of $v,w$ at a time;
conditional-expectation contraction and measure preservation bound the
error uniformly in $u$.  The pair criterion follows for all bounded
functions on $\mathcal R$.

We next prove the following vector assertion on either $X$ or
$\mathcal R$.  On either system use $P$ for its expectation onto $Y$.
If $n_1,\ldots,n_l$ are distinct \emph{nonzero} integers, the $v_j$ are
bounded, and $Pv_{j_*}=0$ for some $j_*$, then
\begin{equation}\label{eq:rigidity-centered-vector}
 \left\|\operatorname{Avg}_{u\in\Phi_N}
          \prod_{j=1}^lU_{u^{n_j}}v_j\right\|_2\longrightarrow0
\end{equation}
for every F\o lner sequence.  Both systems are ergodic and satisfy the
pair criterion, so the same induction proves the assertion on each.
By rescaling we may assume $\norm{v_j}_\infty\leq1$.

For $l=1$, the assertion is mean ergodicity and
$\int v_1=\int Pv_1=0$.  Suppose it is proved for $l-1$, and set
\[
 V_u=\prod_{j=1}^lU_{u^{n_j}}v_j,
 \qquad F_{j,h}=U_{h^{n_j}}v_j\,\bar v_j\quad(h\in K).
\]
For each fixed $h$, invariance of the integral gives
\begin{align}
 \int V_{hu}\overline{V_u}\,d\mu
 &=\int\prod_{j=1}^lU_{u^{n_j}}F_{j,h}\,d\mu\notag\\
 &=\int F_{l,h}\prod_{j<l}U_{u^{n_j-n_l}}F_{j,h}\,d\mu.
 \label{eq:rigidity-vdc-correlation}
\end{align}
The $l-1$ shifted exponents in the last expression are distinct and
nonzero.  Expand each shifted $F_{j,h}$ into $PF_{j,h}$ and its centered
remainder.  The induction hypothesis makes the average of every term
with a centered shifted factor tend to zero in $L^2$; multiplying by the
fixed bounded $F_{l,h}$ does not affect that conclusion after integration.
The remaining shifted product is $Y$-measurable, so the unshifted
$F_{l,h}$ may also be replaced by $PF_{l,h}$ inside its integral.
All the conditional factors are bounded by one.  Hence, for every
F\o lner sequence,
\begin{equation}\label{eq:rigidity-correlation-bound}
 \limsup_{N\to\infty}
 \left|\operatorname{Avg}_{u\in\Phi_N}
       \int V_{hu}\overline{V_u}\,d\mu\right|
 \leq a(h),
 \qquad a(h)=\norm{PF_{j_*,h}}_2.
\end{equation}
This bound also covers $j_*=l$: bound that unshifted factor in $L^1$,
then in $L^2$, and the others in $L^\infty$.  If $j_*<l$, use the same
bound on its shifted factor and measure preservation.  The centered
factor is therefore not lost by moving one exponent to zero.

The pair criterion with $v=\bar v_{j_*}$, $w=v_{j_*}$, and
$d=n_{j_*}\ne0$ shows that averages of $a(h)^2$ tend to zero along
\emph{every} F\o lner sequence; the same is true of $a(h)$ by
Cauchy--Schwarz.  If $(\Psi_M)$ is any chosen F\o lner sequence, it follows
that
\begin{equation}\label{eq:rigidity-uniform-translates}
 \sup_{t\in K}\operatorname{Avg}_{h\in\Psi_M}a(ht)
 \longrightarrow0.
\end{equation}
Indeed, a subsequence of bad translates would itself be a F\o lner
sequence, because $K$ is abelian, contradicting the every-sequence
vanishing just proved.

To finish the induction, fix a finite nonempty shift set $S\subset K$.
The difference between the averages of $V_u$ and of
$|S|^{-1}\sum_{h\in S}V_{hu}$ tends to zero in norm by the F\o lner
property.  Jensen's inequality and expansion in pairs of shifts give
\begin{align*}
 \limsup_{N\to\infty}
 \left\|\operatorname{Avg}_{u\in\Phi_N}V_u\right\|_2^2
 &\leq\frac1{|S|^2}\sum_{h,h'\in S}
 \limsup_{N\to\infty}
 \left|\operatorname{Avg}_{u\in\Phi_N}
          \int V_{hu}\overline{V_{h'u}}\,d\mu\right|\\
 &\leq\frac1{|S|^2}\sum_{h,h'\in S}a(h(h')^{-1}).
\end{align*}
For the second inequality substitute $u'=h'u$ and apply
\eqref{eq:rigidity-correlation-bound} along the translated F\o lner
sequence.  Now take $S=\Psi_M$.  The last expression is bounded by
$\sup_t\operatorname{Avg}_{h\in\Psi_M}a(ht)$ and tends to zero by
\eqref{eq:rigidity-uniform-translates}.  The limits are taken first in
$N$ for fixed finite $S$, and only then in $M$.  This proves
\eqref{eq:rigidity-centered-vector}.

Finally, expand the product in \eqref{eq:source-6} by writing each
$v_j=Pv_j+(v_j-Pv_j)$.  Consider a term
$W_u=\prod_jU_{u^{n_j}}w_j$ with at least one $Pw_j=0$.  Conditional
independence gives
\begin{equation}\label{eq:rigidity-multiple-tensor}
 \norm{PW_u}_2^2
 =\int_{\mathcal R}\prod_j
       U_{u^{n_j}}(w_j\otimes\bar w_j)\,d\mu_{\mathcal R}.
\end{equation}
The corresponding tensor is centered over $Y$, since
$P_{\mathcal R}(w_j\otimes\bar w_j)=|Pw_j|^2$.  In this scalar integral
we may add one common integer to all exponents, by measure preservation,
and choose it so that none of them is zero.  Apply
\eqref{eq:rigidity-centered-vector} on $\mathcal R$.  The average of
\eqref{eq:rigidity-multiple-tensor} tends to zero.  There are finitely
many centered terms, so the squared norm of their sum has vanishing
average as well.  The term with no centered factor is already
$Y$-measurable and equals the product of projections in
\eqref{eq:source-6}.  This proves the lemma.
\end{proof}

\subsection{The radial difference inequality}

Fix an ergodic wild $K$-invariant probability $\mu$ on $D$.  Apply
Lemma~\ref{lem:compact-tower} to the system $(D,\mu)$, and let $Y$ be
its terminal factor.  Lemma~\ref{lem:relative-singularity}, applied to
the identity map of $D$, says that the conditional difference law
$\nu_Y$ over this factor is wild.  Set
\begin{equation}\label{eq:rigidity-projected-characters}
 P=\mathbb E(\,\cdot\mid Y),\qquad
 g_z=P\psi_z,\qquad m(z)=\norm{g_z}_2^2\quad(z\in E).
\end{equation}
Conditional independence identifies its Fourier transform as
\begin{equation}\label{eq:rigidity-difference-fourier}
 \int_D\psi_z\,d\nu_Y=m(z).
\end{equation}
In particular $0\leq m(z)\leq1$, $m(0)=1$, and
Lemma~\ref{lem:line-average} gives \eqref{eq:source-2} with
$f_{\nu_Y}=m$.

Conditional expectation commutes with the action, so
$g_{uz}=U_ug_z$ and $m(uz)=m(z)$ for $u\in K$.  Equal nonzero moduli
in $E$ differ by an element of $K$.  Thus $m$ is radial; its radii lie
in $F_{\geq0}$, and we write $m(r)$ on that set.  Also
$g_{-z}=\bar g_z$ and $\norm{g_z}_\infty\leq1$.

\begin{lemma}[Radial identities]\label{lem:rigidity-radial}
For $w(r)=-\log m(r)\in[0,\infty]$, with $-\log0=\infty$, one has
\begin{equation}\label{eq:source-8}
 w(|A-B|)\leq w(A)+w(B)\qquad(A,B\in F_{\geq0}).
\end{equation}
For every $r\in F$ with $r>0$ and every F\o lner sequence in $K$,
\begin{equation}\label{eq:source-9}
 \operatorname{Avg}_{u\in\Phi_N}m(2r|\operatorname{Re}u|)
 \longrightarrow m(r)^2.
\end{equation}
\end{lemma}

\begin{proof}
First take $A,B>0$ and put $L=\sqrt{AB}\in F$.  For every $u\in K$,
\begin{equation}\label{eq:rigidity-equal-moduli}
 |A-Bu^2|^2
 =(A-B)^2+4AB(\operatorname{Im}u)^2
 =|-L+(A-B)u+Lu^2|^2.
\end{equation}
For the second equality, factor the last expression as
$u((A-B)+L(u-u^{-1}))$ and use $u-u^{-1}=2i\operatorname{Im}u$.

On the first side,
$\psi_{A-Bu^2}=\psi_AU_{u^2}\psi_{-B}$.  By
\eqref{eq:source-6} the projection of this product may be replaced,
in averaged squared $L^2$ error, by $g_AU_{u^2}g_{-B}$.  Both have
$L^2$ norm at most one, so their squared norms differ in averaged
absolute value by a quantity tending to zero.  Indeed the absolute
difference is at most twice their $L^2$ distance, and then
Cauchy--Schwarz applies to the finite average.  Ergodicity of $Y$ and
surjectivity of the square map now give
\begin{align}
 \operatorname{Avg}_{u\in\Phi_N}m(A-Bu^2)
 &\longrightarrow
 \left(\int_Y|g_A|^2\,d\nu\right)
 \left(\int_Y|g_{-B}|^2\,d\nu\right)\notag\\
 &=m(A)m(B).
 \label{eq:rigidity-two-term-limit}
\end{align}
Here the intervening squared norm is
$\int_Y|g_A|^2U_{u^2}|g_{-B}|^2\,d\nu$.

On the second side of \eqref{eq:rigidity-equal-moduli}, apply
\eqref{eq:source-6} with exponents $0,1,2$.  The projected character is
approximated by
\[
 g_{-L}\,U_ug_{A-B}\,U_{u^2}g_L.
\]
Its squared norm is at most $\norm{g_{A-B}}_2^2=m(|A-B|)$, since the
two outside factors have modulus at most one.  Equal moduli give equal
$m$-values, so \eqref{eq:rigidity-two-term-limit} implies
$m(A)m(B)\leq m(|A-B|)$.  Taking negative logarithms proves
\eqref{eq:source-8}, including zero values of $m$ with the stated
extended-real convention.  If $A=0$ or $B=0$, the inequality is an
identity because $w(0)=0$.

For \eqref{eq:source-9}, use
$ru+r\bar u=2r\operatorname{Re}u$ and the distinct exponents $1,-1$.
The same approximation reduces the average to squared norms of
$U_ug_r\,U_{u^{-1}}g_r$.  Those squared norms equal
\[
 \int_Y|g_r|^2U_{u^{-2}}|g_r|^2\,d\nu.
\]
Mean ergodicity on $Y$ through the surjective power map $u\mapsto u^{-2}$
gives the limit $m(r)^2$.  Radiality replaces the real argument by its
absolute value, proving \eqref{eq:source-9}.
\end{proof}

\subsection{General-position rotations and bounded residues}

The remaining task is to show that $m(r)=0$ for every positive $r\in F$.
Assuming instead that $m(r)>0$, we will obtain a uniform lower bound
$m(s)\ge c>0$ on some interval $(0,\varepsilon)\cap F$, contradicting
the vanishing line average.  By \eqref{eq:source-9}, a suitable set
of rotations with bounded $w(2r|\operatorname{Re}u|)$ has finitely many
translates covering $K$.  The first lemma below provides rotations in
rational general position from which the proof will extract bases.
The second controls the resulting integer combinations of nonnegative
lengths under a fixed cap, using only the difference inequality
\eqref{eq:source-8}.

\begin{lemma}[General position in a number field]
\label{lem:rigidity-general-position}
Let $N_0\subset F$ be a number field, let $N=N_0(i)$, and let
$d_0=[N:\Q]$.  There are arbitrarily large finite sets $S\subset K\cap N$
such that every $d_0$ distinct elements of $S$ form a $\Q$-basis of $N$.
\end{lemma}

\begin{proof}
For $t\in N_0$ put
\[
 u(t)=\frac{1+it}{1-it}\in K\cap N.
\]
Choose a $\Q$-basis of $N_0$.  In the rational coordinates of $t$, the
coordinates of $u(t)$ in a $\Q$-basis of $N$ are rational functions over
$\Q$.  This follows, for example, by inverting the matrix for
multiplication by $1-it$.  Its denominator is nonzero for all parameters
in $N_0$, since $t$ is real.

No nonzero $\Q$-linear functional $\ell:N\to\Q$ vanishes on all $u(t)$.
Suppose one did.  For every fixed $a\in N_0$, the rational function
$\ell(u(ha))$ vanishes at every rational $h$, and hence is identically
zero.  Its formal expansion at $h=0$ comes from
\[
 u(ha)=1+2iha-2h^2a^2+\cdots.
\]
The first three coefficients imply
$\ell(1)=\ell(ia)=\ell(a^2)=0$ for every $a\in N_0$.  Squares span
$N_0$ over $\Q$, as
$4a=(a+1)^2-(a-1)^2$.  Thus $\ell$ vanishes on $N_0+iN_0=N$, a
contradiction.  This is a calculation in rational function coordinates
and formal series; it uses no usual-topology continuity of $\ell$.

For any proper rational subspace of $N$, choose a nonzero rational
linear functional annihilating it.  Its value on $u(t)$ has a nonzero
numerator polynomial in the rational coordinates of $t$.  A finite
union of such proper subspaces therefore cannot contain all $u(t)$:
the product of their nonzero numerator polynomials is nonzero, and a
nonzero polynomial over the infinite field $\Q$ cannot vanish on all
of rational affine space.  Inductively choose each new element outside
the spans of all subsets of at most $d_0-1$ earlier elements.  Every
subset of at most $d_0$ chosen elements remains independent.  This
constructs finite sets of every required size.
\end{proof}

\begin{lemma}[Addition under a cap]\label{lem:rigidity-capped-residues}
Let $w:F_{\geq0}\to[0,\infty]$ satisfy $w(0)=0$ and
\eqref{eq:source-8}.  Let $p_1,\ldots,p_q\in F_{\geq0}$ with
$p=\max_jp_j>0$ and $w(p_j)\leq H<\infty$.  If
\[
 z=\sum_{j=1}^q a_jp_j\in[0,p),\qquad
 a_j\in\Z,\qquad\sum_j|a_j|\leq M,
\]
then $w(z)\leq4MH$.
\end{lemma}

\begin{proof}
Keep every partial sum as its residue in $[0,p)$.  A positive input
$p_j<p$ has cost at most $H$, while the input $p$ has zero residue.
The inverse of a nonzero residue $b$ is $p-b$ and satisfies
$w(p-b)\leq w(p)+w(b)$.  Thus each signed input $\pm p_j$ has a
residue of cost at most $2H$; the zero cases cause no difficulty.

For two current residues $a,b\in[0,p)$, if $a+b<p$, use three successive
applications of the difference inequality:
\begin{align*}
 w(a+b)
 &=w\bigl(p-((p-a)-b)\bigr)\\
 &\leq w(p)+w((p-a)-b)\\
 &\leq w(p)+w(p-a)+w(b)\\
 &\leq2w(p)+w(a)+w(b).
\end{align*}
All arguments are nonnegative.  If $a+b>p$, the new residue is
$a-(p-b)$, whose cost is at most $w(a)+w(p)+w(b)$.  If $a+b=p$, its
cost is zero.  Adding a signed input therefore increases the current
cost by at most $4H$.  Expand the integer combination into at most $M$
signed inputs and start from residue zero.  Its final residue is $z$
itself, proving the bound.
\end{proof}

\subsection{Completion of the rigidity theorem}

\begin{proof}[Proof of Theorem~\ref{thm:rigidity}]
First let $\mu$ be ergodic and wild, and use the factor and radial
function constructed above.  Suppose that $m(r)>0$ for some $r\in F$,
$r>0$.  Choose $H>0$ finite so that $e^{-H}<m(r)^2$.  The set
\[
 A_0=\{u\in K:w(2r|\operatorname{Re}u|)\leq H\}
\]
is syndetic.  Indeed, if no finite translates cover $K$, then for every
finite $S\subset K$ some translate $Sv$ avoids $A_0$.  Apply this to
the sets of a F\o lner sequence and choose the avoiding translate
separately for each set.  Those translates still form a F\o lner
sequence, whereas all their averaged $m$-values are at most $e^{-H}$.
This contradicts \eqref{eq:source-9}.  Hence there are
$q_1,\ldots,q_J\in K$ such that
\begin{equation}\label{eq:rigidity-syndetic-cover}
 \text{for every }v\in K\text{, some }q_jv\text{ belongs to }A_0.
\end{equation}

Choose a number field $N_0\subset F$ containing the real and imaginary
parts of every $q_j$, and put $N=N_0(i)$ and $d_0=[N:\Q]$.
Lemma~\ref{lem:rigidity-general-position} supplies a finite
$S\subset K\cap N$ of size greater than $J(d_0-1)$ in rational linear
general position.  For each $v\in K$, apply
\eqref{eq:rigidity-syndetic-cover} to all the points $sv$, $s\in S$.
Pigeonhole gives an index $j$ and $d_0$ elements $s$ for which
$q_jsv\in A_0$.  Their coefficients $k=q_js$ form a $\Q$-basis of $N$,
since multiplication by $q_j$ is an invertible $\Q$-linear map.
Call such a basis a good selection for $v$.

There are only finitely many possible selections.  Clearing the rational
denominators in the basis expression of $1$, choose one positive integer
$D_0$ such that, for each possible selection $\mathcal T$,
\begin{equation}\label{eq:rigidity-integer-selection}
 D_0\,1=\sum_{k\in\mathcal T}a_k k,
 \qquad a_k\in\Z,\qquad
 \sum_{k\in\mathcal T}|a_k|\leq M_0,
\end{equation}
where the finite bound $M_0$ is common to all selections.  There is also
a common $c_0>0$ such that
\begin{equation}\label{eq:rigidity-selection-transversality}
 \max_{k\in\mathcal T}|\operatorname{Re}(kv)|\geq c_0
 \qquad(v\in\T)
\end{equation}
for every selection.  To prove this, for a fixed selection the maximum
is a continuous function of the ordinary unit direction $v$.  It never
vanishes: the rational span of the selection contains $1$ and $i$, so
vanishing would imply both $\operatorname{Re}v=0$ and
$\operatorname{Re}(iv)=0$.  Take a positive minimum on the compact unit
circle and then the minimum over the finitely many selections.  We may
and do take $c_0\leq1$.

Now let $v\in K$ satisfy $|\operatorname{Re}v|<c_0/D_0$, and choose a
good selection $\mathcal T$ for it.  Its nonnegative numbers
\[
 p_k=2r|\operatorname{Re}(kv)|\quad(k\in\mathcal T),
 \qquad p=\max_{k\in\mathcal T}p_k
\]
satisfy $w(p_k)\leq H$ and $p\geq2rc_0$.  Taking real parts of
\eqref{eq:rigidity-integer-selection}, and inserting the signs of the
real parts and then the overall sign, expresses
\[
 z=2rD_0|\operatorname{Re}v|<p
\]
as an integer combination of the $p_k$ with coefficient length at most
$M_0$.  All these numbers lie in $F$.  Lemma~\ref{lem:rigidity-capped-residues}
therefore gives
\begin{equation}\label{eq:rigidity-small-radius-bound}
 w(2rD_0|\operatorname{Re}v|)\leq4M_0H
 \qquad\left(v\in K,\ |\operatorname{Re}v|<c_0/D_0\right).
\end{equation}

Every $t\in F\cap[-1,1]$ is the real part of an element of $K$, namely
$t+i\sqrt{1-t^2}$.  For every $s\in F\cap(0,2rc_0)$ apply
\eqref{eq:rigidity-small-radius-bound} with
$t=s/(2rD_0)$.  It follows that
\[
 m(s)\geq e^{-4M_0H}>0
 \qquad\bigl(s\in F\cap(0,2rc_0)\bigr).
\]
Choose a nonnegative $\phi\in C_c(\R)$ supported in this open interval
with positive integral.  Positivity of the functional $I$ and its
agreement with Lebesgue integration on continuous test functions give
\[
 I\bigl(\phi(t)m(t)\bigr)
 \geq e^{-4M_0H}I(\phi)
 =e^{-4M_0H}\int_\R\phi(t)\,dt>0.
\]
This contradicts \eqref{eq:source-2} for the wild law $\nu_Y$, with
$a=1$, in view of \eqref{eq:rigidity-difference-fourier}.  This step
uses no ordinary-measure measurability of the radial function: $I$ is
defined on all the bounded, boundedly supported functions in question.

Thus $m(r)=0$ for every positive $r\in F$, and radiality gives $m(z)=0$
for all nonzero $z\in E$.  Consequently
\[
 \left|\int_D\psi_z\,d\mu\right|
 =\left|\int_Yg_z\,d\nu\right|
 \leq\norm{g_z}_2=0\qquad(z\ne0).
\]
Fourier uniqueness on the compact abelian group $D$ identifies $\mu$
with Haar probability.

Finally, an arbitrary wild invariant probability has an ergodic
decomposition for the countable group $K$ on the standard compact space
$D$.  Since $C$ is invariant Borel and has measure zero, almost every
ergodic component also gives $C$ measure zero.  The ergodic case just
proved makes each such component Haar probability.  Their integral is
therefore Haar probability as well.  This proves the theorem.
\end{proof}
